%E1===============================[=.............................]===============================%E1
\def\EoneEN{Which square is cut into two different shapes?}
\def\EoneLT{Kuris kvadratas padalytas į dvi skirtingas figūras?}
\def\EonePL{Który kwadrat został rozcięty wzdłuż pogrubionych linii
na dwie części o różnych kształtach?}
\def\EoneRU{Какой квадрат разрезан на две разные фигуры?}

\def\rEone{E}
\def\aEone{\granswer[0.7]{E01-1}{E01-2}{E01-3}{E01-4}{E01-5}}

%E2===============================[==............................]===============================%E2
\def\EtwoEN{Which of the following sums is the largest?}
\def\EtwoLT{Kuri iš šių sumų yra didžiausia?}
\def\EtwoPL{Która z następujących sum jest największa?}
\def\EtwoRU{Которая из перечисленных сумм наибольшая?}

\def\rEtwo{A}
\def\aEtwo{\answer{202+4}{20+24}{2+0+2+4}{20+2+4}{2+0+24}}

%E3===============================[===...........................]===============================%E3
\def\EthreeEN{The table consists of 28 white cells. Ira paints two rows and one column.  A row is from left to right.  A column is from top to bottom.  How many cells will remain white?}
\def\EthreeLT{Lentelėje yra 28 langeliai. Jorė nuspalvino dvi eilutes ir vieną stulpelį. Kiek liko neneuspalvintų langelių?}
\def\EthreePL{Tabelka składa się z \mat{28} pól. Dominika pokolorowała wszystkie pola znajdujące się w dwóch wierszach i wszystkie pola w jednej kolumnie tej tabelki. Ile pól pozostało niezamalowanych?}
\def\EthreeRU{Таблица состоит из 28 белых клеток (см. рисунок).
Ирина закрашивает две строки и один столбец.
Сколько клеток останутся белыми?}

\def\rEthree{C}
\def\aEthree{\answer{8}{10}{12}{14}{16}}

%E4===============================[====..........................]===============================%E4
\def\EfourEN{Soccer players numbered 1 to 11 stand in a circle.  Each player kicks the ball to the third player on their left.
Player 1 starts.  This kicking pattern continues until a player has the ball for the second time.  What is the number of the player who kicked the ball last?}
\def\EfourLT{Futbolininkai, sunumeruoti nuo 1 iki 11, stovi ratu veidu į centrą. Kiekvienas žaidėjas perduoda kamuolį
trečiam į kairę nuo jo žaidėjui (žr. paveikslėlį). Pradeda žaidėjas numeris 1. Perdavimai tęsiasi tol,
kol kamuolys pas kurį nors iš žaidėjų atsiduria antrą kartą. Koks numeris žaidėjo,
    kuris perdavė kamuolį paskutinis?}
\def\EfourPL{Piłkarze z numerami od \mat{1} do \mat{11} stoją w kręgu przodem do siebie. Gracz z~numerem \mat{1} ma piłkę i~podaje ją do trzeciego gracza po swojej lewej stronie. Ten, po otrzymaniu piłki, też podaje ją do trzeciego gracza po swojej lewej stronie (patrz rysunek).
Taki schemat podań powtarza się, dopóki
piłka nie trafi do któregoś z graczy
po raz drugi. Jaki jest numer zawodnika, który jako ostatni podał piłkę?}
\def\EfourRU{Футболисты с номерами от 1 до 11 встают в круг лицом к центру.
Каждый игрок передаёт мяч третьему игроку слева от себя (см. рисунок). Начинает игрок с номером 1.
Такие передачи продолжаются до тех пор, пока один из игроков не получит мяч во второй раз. Какой номер у игрока, передавшего мяч последним?}

\def\rEfour{B}
\def\aEfour{\answer{11}{9}{8}{6}{4}}

%E5===============================[=====.........................]===============================%E5
\def\EfiveEN{Mohammad wrote three consecutive 4-digit numbers in a row.  Then he hid some digits (see picture).  What are the hidden digits?}
\def\EfiveLT{Matas užrašė tris iš eilės einančius keturženklius skaičius. Tada kai kuriuos skaitmenis jis
uždengė (žr. paveikslėlį). Kas buvo uždengta?}
\def\EfivePL{Staś zapisał trzy kolejne liczby czterocyfrowe. Następnie zasłonił niektóre cyfry (patrz rysunek). Co zostało zakryte?}
\def\EfiveRU{Марк написал три последовательных четырёхзначных числа и закрыл некоторые цифры (см. рисунок).
Какие цифры он закрыл?}

\def\rEfive{D}
\def\aEfive{\answer{389, 3, 99}{489, 3, 96}{489, 4, 98}{489, 4, 99}{488, 4, 99}}

%E6===============================[======........................]===============================%E6
\def\EsixEN{Lizzy pays 7 dollars for three items.  The cost of each item is different and is a whole number.  How much is the most expensive item?}
\def\EsixLT{Už tris skirtingus pyragaičius Adelė užmokėjo 7 eurus. Kiekvienas pyragaitis kainavo sveikąjį skaičių
eurų, visos trys kainos skyrėsi. Kiek kainavo brangiausias pyragaitis?}
\def\EsixPL{Za trzy różne ciastka Helenka zapłaciła \mat{7} złotych. Każde ciastko kosztowało całkowitą liczbę złotych i
miało inną cenę. Ile kosztowało najdroższe ciastko?}
\def\EsixRU{Элина платит 7 евро за три пирожных. Цена каждого пирожного является целым числом и
отличается от цены каждого другого пирожного. Сколько стоит самое дорогое пирожное?}

\def\rEsix{C}
\def\aEsixEN{\answer{2 dollars}{3 dollars}{4 dollars}{5 dollars}{6 dollars}}
\def\aEsixLT{\answer{2 eurus}{3 eurus}{4 eurus}{5 eurus}{6 eurus}}
\def\aEsixPL{\answer{2 złote}{3 złote}{4 złote}{5 złotych}{6 złotych}}
\def\aEsixRU{\answer{2 евро}{3 евро}{4 евро}{5 евро}{6 евро}}

%E7===============================[=======.......................]===============================%E7
\def\EsevenEN{A cat knocks off one block from Felix's construction. What could this construction have looked like before the block was knocked off?}
\def\EsevenLT{Gediminas iš trinkelių sustatė statinį. Katinas iš to statinio nuspyrė vieną trinkelę (žr.
    pavеikslėlį). Kaip galėjo atrodyti tas statinys dar su visomis trinkelėmis?}
\def\EsevenPL{Ignaś zbudował konstrukcję z klocków. Kot strącił z niej jeden klocek --- patrz obrazek. Jak mogła wyglądać ta budowla przed strąceniem klocka?}
\def\EsevenRU{Кот сбивает с конструкции Феликса один кубик (см. рисунок). Как могла выглядеть эта конструкция до того, как кот сбил кубик?}

\def\rEseven{E}
\def\aEseven{\granswer[0.5]{E07-2}{E07-3}{E07-4}{E07-5}{E07-6}}

%E8===============================[========......................]===============================%E8
\def\EeightEN{Alex has a Kangaroo poster on the kitchen wall. How many grey tiles are there behind the poster?}
\def\EeightLT{Alius virtuvėje ant sienos prisiklijavo plakatą „Kengūra 2024“.
Kelias tamsiąsias plyteles uždengė plakatas?
}
\def\EeightPL{Ściana kuchni była wyłożona naprzemiennie białymi i szarymi płytkami. Marysia nakleiła na niej prostokątny plakat konkursu Kangur --- patrz obrazek. Ile szarych płytek przykrył ten plakat?}
\def\EeightRU{У Алексея на стене кухни висит плакат «Кенгуру 2024». Сколько серых плиток находится за плакатом?}

\def\rEeight{B}
\def\aEeight{\answer{15}{21}{25}{30}{35}}

%E9===============================[=========.....................]===============================%E9
\def\EnineEN{Lucian wants to choose two of the six numbers 1, 2, 3, 4, 5, 6 and find their sum.
How many different results can he get?}
\def\EnineLT{Elena turi pasirinkti du iš šešių skaičių 1, 2, 3, 4, 5, 6 ir rasti jų sumą. Kiek skirtingų
rezultatų ji gali gauti?}
\def\EninePL{Jakub chce wybrać dwie z sześciu  liczb 1, 2, 3, 4, 5, 6 i obliczyć ich sumę.
Ile różnych rezultatów może w ten sposób otrzymać?}
\def\EnineRU{Стасик хочет выбрать два из шести чисел 1, 2, 3, 4, 5, 6 и подсчитать их сумму. Сколько
различных результатов он может получить?}

\def\rEnine{C}
\def\aEnine{\answer{7}{8}{9}{15}{30}}

%E10===============================[==========....................]===============================%E10
\def\EtenEN{There are five fruits in a bowl:
\includegraphics[scale=0.7]{E10-1} \includegraphics[scale=0.7]{E10-2} \includegraphics[scale=0.7]{E10-3} \includegraphics[scale=0.7]{E10-4} \includegraphics[scale=0.7]{E10-5}.
Ann likes \includegraphics[scale=0.7]{E10-3}~\includegraphics[scale=0.7]{E10-4}~\includegraphics[scale=0.7]{E10-5}.
Ben likes \includegraphics[scale=0.7]{E10-3}.
Cam likes \includegraphics[scale=0.7]{E10-1} \includegraphics[scale=0.7]{E10-2} \includegraphics[scale=0.7]{E10-3} \includegraphics[scale=0.7]{E10-4} \includegraphics[scale=0.7]{E10-5}.
Dan likes \includegraphics[scale=0.7]{E10-3} \includegraphics[scale=0.7]{E10-5}.
Eli likes \includegraphics[scale=0.7]{E10-2} \includegraphics[scale=0.7]{E10-4}.
\newline Everyone gets a fruit they like. What does Ben get?}
\def\EtenLT{Ant padėklo buvo penki vaisiai:\includegraphics[scale=0.7]{E10-1} \includegraphics[scale=0.7]{E10-2} \includegraphics[scale=0.7]{E10-3} \includegraphics[scale=0.7]{E10-4} \includegraphics[scale=0.7]{E10-5}.
Emilija mėgsta \includegraphics[scale=0.7]{E10-3} \includegraphics[scale=0.7]{E10-4} \includegraphics[scale=0.7]{E10-5},
Jonas mėgsta \includegraphics[scale=0.7]{E10-3},
Lukas mėgsta \includegraphics[scale=0.7]{E10-1} \includegraphics[scale=0.7]{E10-2} \includegraphics[scale=0.7]{E10-3} \includegraphics[scale=0.7]{E10-4} \includegraphics[scale=0.7]{E10-5},
Agnė mėgsta \includegraphics[scale=0.7]{E10-3} \includegraphics[scale=0.7]{E10-5},
Liepa mėgsta \includegraphics[scale=0.7]{E10-2}~\includegraphics[scale=0.7]{E10-4}.
Kiekvienas vaikas gavo po mėgiamą vaisių. Ką gavo Lukas?}
\def\EtenPL{Na tacy leżało pięć owoców: \includegraphics[scale=0.7]{E10-1} \includegraphics[scale=0.7]{E10-2} \includegraphics[scale=0.7]{E10-3} \includegraphics[scale=0.7]{E10-4} \includegraphics[scale=0.7]{E10-5}.
%
Kajtek lubi \includegraphics[scale=0.7]{E10-3} \includegraphics[scale=0.7]{E10-4} \includegraphics[scale=0.7]{E10-5},
Franek lubi~\includegraphics[scale=0.7]{E10-3},
Józek lubi \includegraphics[scale=0.7]{E10-1} \includegraphics[scale=0.7]{E10-2} \includegraphics[scale=0.7]{E10-3} \includegraphics[scale=0.7]{E10-4} \includegraphics[scale=0.7]{E10-5},
Tadzio lubi  \includegraphics[scale=0.7]{E10-3} \includegraphics[scale=0.7]{E10-5},
Alicja lubi \includegraphics[scale=0.7]{E10-2} \includegraphics[scale=0.7]{E10-4}. Każde dziecko dostało  taki owoc, jaki lubi. Co dostał Józek?}
\def\EtenRU{В вазе было пять фруктов: \includegraphics[scale=0.7]{E10-1} \includegraphics[scale=0.7]{E10-2} \includegraphics[scale=0.7]{E10-3} \includegraphics[scale=0.7]{E10-4} \includegraphics[scale=0.7]{E10-5}.
Аня любит \includegraphics[scale=0.7]{E10-3} \includegraphics[scale=0.7]{E10-4} \includegraphics[scale=0.7]{E10-5}.
Борис любит~\includegraphics[scale=0.7]{E10-3}.
Саша любит \includegraphics[scale=0.7]{E10-1} \includegraphics[scale=0.7]{E10-2} \includegraphics[scale=0.7]{E10-3} \includegraphics[scale=0.7]{E10-4} \includegraphics[scale=0.7]{E10-5}.
Дима любит \includegraphics[scale=0.7]{E10-3} \includegraphics[scale=0.7]{E10-5}.
Ева любит \includegraphics[scale=0.7]{E10-2}~\includegraphics[scale=0.7]{E10-4}.
Каждый из детей получил тот фрукт, который он любит. Что получил Саша?}

\def\rEten{A}
\def\aEten{\gransweranswer[0.7]{E10-1}{E10-2}{E10-3}{E10-4}{E10-5}}

%E11===============================[===========...................]===============================%E11
\def\EuoneEN{Ada has built a tower of 8 discs, as in the picture. Ada removes the second disc from the bottom of this tower.  Then she removes the third disc from the bottom of the new tower.  Then she removes the fourth disc from the bottom of the new tower. Then she removes the fifth disc from the bottom of the new tower.  Which tower does Ada end up with?}
\def\EuoneLT{Kristina iš 8 skridinių sustatė bokštą (žr. paveikslėlį). Iš to bokšto ji išėmė antrą nuo apačios
skridinį. Tada iš susidariusio bokšto ji išėmė trečią nuo apačios skridinį. Iš naujai susidariusio bokšto ji
pašalino ketvirtą nuo apačios skridinį.  Pagaliau iš naujojo bokšto ji pašalino penktą nuo apačios skridinį.
Koks susidarė bokštas?}
\def\EuonePL{Ada zbudowała wieżę z ośmiu żetonów --- patrz obrazek.\\
Usunęła z niej drugi żeton od dołu.\\
Później usunęła trzeci od dołu  żeton  nowej wieży.\\
Z otrzymanej wieży usunęła  żeton czwarty od dołu.\\
Z tak powstałej wieży usunęła żeton piąty od dołu.
Jaką wieżę Ada otrzymała na koniec?}
\def\EuoneRU{Ада построила башню из 8 дисков (см. рисунок). Из этой башни Ада убирает второй снизу диск.
Затем из новой башни она убирает третий снизу диск.
После этого из получившейся башни она убирает четвёртый снизу диск.
И, наконец, она убирает пятый снизу диск новой башни. Какая башня получилась у Ады?}

\def\rEuone{B}
\def\aEuone{\granswer[0.8]{E11-2}{E11-3}{E11-4}{E11-5}{E11-6}}

%E12===============================[============..................]===============================%E12
\def\EutwoEN{Peter the penguin goes fishing every day and brings back 9 fish for his 2 chicks. Each day, he gives 5 fish to the first chick he sees and 4 fish to the second chick, which they eat. Over the last few days, 1 chick has eaten 26 fish. How many fish has the other chick eaten?}
\def\EutwoLT{Pingvinas Ginas kasdien žvejoja ir parsineša dviem savo cypliams 9 žuvis. Kiekvieną dieną
pirmam pamatytam cypliui atiduoda 5 žuvis, o antram -- 4 žuvis. Per kelias paskutines dienas vienas iš cyplių sušveitė
26 žuvis. Kiek žuvų per tą laiką atiteko kitam cypliui?}
\def\EutwoPL{Pingwin Czarnuś codziennie przynosi \mat{9} ryb swoim dwojgu pisklętom. Każdego dnia daje pierwszemu napotkanemu pisklęciu \mat{5}~ryb, a drugiemu \mat{4} ryby. W~ciągu ostatnich kilku dni jedno z piskląt otrzymało \mat{26} ryb. Ile ryb dostało w tym czasie drugie pisklę?}
\def\EutwoRU{Пингвин каждый день ходит на рыбалку и приносит 9 рыб для своих двух птенцов.
Каждый день он даёт первому увиденному птенцу 5 рыб, а второму – 4 рыбы, которые они съедают.
За последние несколько дней один из птенцов съел 26 рыб. Сколько рыб съел другой птенец?}

\def\rEutwo{D}
\def\aEutwo{\answer{19}{22}{25}{28}{31}}

%E13===============================[=============.................]===============================%E13
\def\EuthreeEN{Seven cards, numbered 1 to 7, are placed in four overlapping rings.  The sum of the numbers in each ring is 10.  Which number is under the question mark?}
\def\EuthreeLT{Septynios kortelės su skaičiais nuo 1 iki 7 padėtos į keturis persidengiančius
vielinius žiedus (žr.~paveikslėlį). Kiekviename žiede kortelėse parašytų skaičių suma lygi 10. Koks skaičius
    parašytas kortelėje, pažymėtoje klaustuku?}
\def\EuthreePL{Siedem kart z liczbami od \mat{1} do \mat{7} ułożono w~czterech  okręgach. Dwie karty odkryto --- patrz rysunek. Suma liczb w każdym okręgu wynosi \mat{10}.
Jaka liczba jest na karcie ze znakiem zapytania?}
\def\EuthreeRU{Семь карточек с числами от 1 до 7 помещены в четыре перекрывающихся кольца из проволоки (см. рисунок). Сумма чисел на карточках в каждом кольце равна 10.
Только две карточки лежат так, что видны числа. Какое число написано на карточке с вопросительным знаком?}

\def\rEuthree{A}
\def\aEuthree{\answer{1}{2}{4}{5}{7}}

%E14===============================[==============................]===============================%E14
\def\EufourEN{ Lucas wants to make a caterpillar that has a head, a tail and either 1, 2 or 3 puzzle pieces in between.
\begin{center}
\includegraphics{E14-1EN}
\end{center}
How many different caterpillars can Lucas make without flipping pieces?}
\def\EufourLT{Julius turi penkias dėlionės detales ir nori iš jų sudėti vikšrą, kuris turėtų galvą,
uodegą ir arba vieną, arba dvi, arba visas tris detales tarp jų (žr. paveikslėlį).
    \begin{center}
\includegraphics{E14-1LT}
\end{center}
Keliais skirtingais būdais
Julius gali sudėti tokį vikšrą? (Detalių sukioti negalima.)}
\def\EufourPL{Julek ma \mat{5} puzzli i chce ułożyć z nich  gąsienicę, która ma głowę, ogon i albo jeden, albo dwa, albo trzy inne elementy układanki pomiędzy nimi.
\begin{center}
\includegraphics{E14-1PL}
\end{center}
Na ile różnych sposobów Julek może zbudować taką gąsienicę? (Puzzli nie wolno odwracać.)}
\def\EufourRU{Лука хочет сделать гусеницу с головой, хвостом и либо одной, либо двумя, либо всеми тремя деталями между ними (см. рисунок).
\begin{center}
\includegraphics{E14-1RU}
\end{center}
Поворачивать или переворачивать детали нельзя. Сколько разных гусениц может построить Лука?}

\def\rEufour{B}
\def\aEufour{\answer{3}{4}{5}{6}{7}}

%E15===============================[===============...............]===============================%E15
\def\EufiveEN{John writes the numbers 1 to 4 on a sheet.  Then he flips the sheet and writes the numbers 5 to 8, as shown.
After that, he cuts the sheet into 4 rectangular cards and puts them in a row:
\begin{center}
    \includegraphics[scale=0.4]{E15-1}
\end{center}
\mbox{What is the sum of the numbers represented by the question marks?}}
\def\EufiveLT{Tamsiojoje lapo pusėje Romas parašė skaičius 1, 2, 3, 4 (žr.~paveikslėlį). Tada per dešinę kraštinę
apvertė lapą šviesiąja puse.
Joje jis parašė skaičius 5, 6, 7, 8. Pagaliau jis sukarpė lapą į keturias lygias dalis:
\vspace{-8pt}
\begin{center}
    \includegraphics[scale=0.4]{E15-1}
\end{center}
\vspace{-5pt}
Kam lygi po klaustukais
paslėptų skaičių suma?}
\def\EufivePL{Na zacieniowanej stronie kartki Ania zapisała cztery liczby.
Następnie, wzdłuż krawędzi przechodzącej przez przerywaną linię,  odwróciła kartkę na białą stronę.
Zapisała na niej cztery inne liczby -- patrz obrazek. Na koniec przecięła kartkę na cztery równe części
\vspace{-10pt}
\begin{center}
    \includegraphics[scale=0.4]{E15-1}
\end{center}
\vspace{-5pt}
Ile wynosi suma liczb ukrytych pod znakami zapytania?}
\def\EufiveRU{На тёмной стороне листа бумаги Рем пишет числа от 1 до~4.
После этого он переворачивает лист на светлую сторону и пишет числа от 5 до 8 (см. рисунок).
Затем он разрезает лист на 4 прямоугольные карточки:
\vspace{-3pt}
\begin{center}
    \includegraphics[scale=0.4]{E15-1}
\end{center}
\vspace{-5pt}
\mbox{Какова сумма чисел, обозначенных вопросительными знаками?}}

\def\rEufive{B}
\def\aEufive{\answer{3}{4}{5}{6}{7}}

%E16===============================[================..............]===============================%E16
\def\EusixEN{A floor is covered with two kinds of tiles \includegraphics[scale=0.9]{E16-1} and~\includegraphics[scale=0.9]{E16-2}.
The rectangles have size $23\ \textup{cm} \times 11\ \textup{cm}$.
The picture shows a part of the floor. What is the side-length of the square tiles?}
\def\EusixLT{Grindys išgrįstos stačiakampėmis \includegraphics{E16-1} ir kvadratinėmis \includegraphics{E16-2} plytelėmis (grindų
    fragmentas pavaizduotas paveikslėlyje). Stačiakampės plytelės matmenys yra
    $23\ \textup{cm} \times 11\ \textup{cm}$. Koks yra kvadratinės plytelės kraštinės ilgis?}
\def\EusixPL{Podłoga jest wyłożona płytkami dwóch rodzajów: jasnymi prostokątnymi \includegraphics{E16-1} i ciemnymi kwadratowymi \includegraphics{E16-2}. Jasna płytka ma wymiary \mat{\unit[23]{cm}\times\unit[11]{cm}}. Fragment tej podłogi jest pokazany na obrazku. Jaką długość ma bok płytki kwadratowej?}
\def\EusixRU{Пол покрыт плитками двух видов:\newline прямоугольной~\includegraphics{E16-1} и квадратной \includegraphics{E16-2}. На рисунке видна часть пола.
Длина большей стороны прямоугольной плитки равна 23~см, меньшей~--~11~см. Какова длина стороны квадратной плитки?}

\def\rEusix{D}
\def\aEusixEN{\answer{$3\ \textup{cm}$}{$4\ \textup{cm}$}{$5\ \textup{cm}$}{$6\ \textup{cm}$}{$7\ \textup{cm}$}}
\def\aEusixLT{\answer{$3\ \textup{cm}$}{$4\ \textup{cm}$}{$5\ \textup{cm}$}{$6\ \textup{cm}$}{$7\ \textup{cm}$}}
\def\aEusixPL{\answer{$3\ \textup{cm}$}{$4\ \textup{cm}$}{$5\ \textup{cm}$}{$6\ \textup{cm}$}{$7\ \textup{cm}$}}
\def\aEusixRU{\answer{$3\ \textup{см}$}{$4\ \textup{см}$}{$5\ \textup{см}$}{$6\ \textup{см}$}{$7\ \textup{см}$}}

%E17===============================[=================.............]===============================%E17
\def\EusevenEN{The sum of three integers is 782. Unfortunately, ink was spilled on three digits (see picture). What is the sum of the three hidden digits?}
\def\EusevenLT{Mindaugas sudėjo tris triženklius skaičius ir užrašė jų sumą 782.
Nelaimei, jis sutepė rašalu tris skaitmenis. Kokia yra neįskaitomųjų skaitmenų suma?}
\def\EusevenPL{Kostek dodał trzy liczby trzycyfrowe i zapisał ich sumę \mat{782}. Niestety, zachlapał atramentem trzy cyfry. Jaka jest suma cyfr pod trzema kleksami?}
\def\EusevenRU{Оля сложила три числа. Их сумма равна~782. Братик Оли
заляпал чернилами три цифры (см. рисунок). Какова сумма трёх заляпанных~цифр?}

\def\rEuseven{D}
\def\aEuseven{\answer{8}{9}{10}{11}{12}}

%E18===============================[==================............]===============================%E18
\def\EueightEN{Lucy weighs three blocks in pairs.
\begin{center}
    \includegraphics[scale=0.3]{E18-1}
\end{center}
How much do all three blocks weigh together?}
\def\EueightLT{Ieva tris skirtingas figūras pasvėrė po dvi, rezultatai parodyti paveikslėlyje.
\begin{center}
    \includegraphics[scale=0.3]{E18-1}
\end{center}
\vspace{-5pt}
Kiek tos trys
figūros sveria kartu?}
\def\EueightPL{Ewa ma trzy różne bryły. Kładła je po dwie na wadze i otrzymała następujące wyniki --- patrz obrazki.
\begin{center}
    \includegraphics[scale=0.3]{E18-1}
\end{center}
\vspace{-5pt}
Ile ważą razem te trzy bryły?}
\def\EueightRU{Аня взвешивает три фигуры по две, как показано на рисунке.
\begin{center}
    \includegraphics[scale=0.3]{E18-1RU}
\end{center}
\vspace{-5pt}
Каков общий вес всех трёх фигур?}

\def\rEueight{A}
\def\aEueightEN{\answer{$270\ \textup{g}$}{$280\ \textup{g}$}{$290\ \textup{g}$}{$300\ \textup{g}$}{$310\ \textup{g}$}}
\def\aEueightLT{\answer{$270\ \textup{g}$}{$280\ \textup{g}$}{$290\ \textup{g}$}{$300\ \textup{g}$}{$310\ \textup{g}$}}
\def\aEueightPL{\answer{$270\ \textup{g}$}{$280\ \textup{g}$}{$290\ \textup{g}$}{$300\ \textup{g}$}{$310\ \textup{g}$}}
\def\aEueightRU{\answer{$270\ \textup{г}$}{$280\ \textup{г}$}{$290\ \textup{г}$}{$300\ \textup{г}$}{$310\ \textup{г}$}}

%E19===============================[===================...........]===============================%E19
\def\EunineEN{There are 60 pupils on a trip. When they line up, the colours of their reflective vests follow the pattern: yellow, green, yellow, green... The colours of their backpacks follow a different pattern: red, brown, orange, red, brown, orange... How many pupils with a yellow reflective vest also have an orange backpack?}
\def\EunineLT{Prieš ekskursiją 60 mokinių išsirikiavo viena eile. Jų apsauginių striukių spalva kartojosi taip:
geltona, žalia, geltona, žalia, ... . Jų kepurių spalva kartojosi kitaip: raudona, ruda, mėlyna, raudona,
ruda, mėlyna, ... . Kelių mokinių su geltona striuke kepurė buvo mėlyna?}
\def\EuninePL{Przed wycieczką \mat{60} dzieci ustawiło się w szeregu. Dwa kolory ich kamizelek odblaskowych,
    zaczynając od pierwszego dziecka, to naprzemiennie: żółty, zielony, żółty, zielony \ldots\
    Trzy kolory plecaków dzieci, rozpoczynając od pierwszego, powtarzają się według schematu:
    czerwony, brązowy, niebieski, czerwony, brązowy, niebieski \ldots\ Ilu uczniów w żółtej kamizelce miało niebieski plecak?}
\def\EunineRU{В поездке участвуют 60 учеников. Когда они выстраиваются в ряд, цвета их светоотражающих жилетов повторяются так:
жёлтый, зелёный, жёлтый, зелёный,~... . Цвета их рюкзаков повторяются по-другому: красный, коричневый, оранжевый, красный, коричневый, оранжевый, ... .
У какого количества учеников жёлтый светоотражающий жилет и оранжевый рюкзак?}

\def\rEunine{E}
\def\aEunine{\answer{3}{4}{6}{8}{10}}

%E20===============================[====================..........]===============================%E20
\def\EutenEN{In the following calculations, the same digits are hidden under the same figures. Different digits are hidden under different figures. What is the value of
\includegraphics[scale=0.5]{E20-2} ?}
\def\EutenLT{Paveikslėlyje vienodos figūros slepia vienodus skaitmenis, skirtingos figūros -- skirtingus skaitmenis.
Kam lygi\newline sandauga \includegraphics[scale=0.5]{E20-2} ?}
\def\EutenPL{Na diagramie obok pod taką samą figurą jest ukryta taka sama cyfra, pod różnymi figurami różne cyfry. Ile wynosi wynik mnożenia \includegraphics[scale=0.5]{E20-2} ?}
\def\EutenRU{В примерах (см. рисунок) одинаковые числа закрыты одинаковыми фигурами, разные числа закрыты разными фигурами.
Чему равно произведение \includegraphics[scale=0.5]{E20-2} ?}

\def\rEuten{C}
\def\aEuten{\answer{0}{15}{18}{28}{30}}

%E21===============================[=====================.........]===============================%E21
\def\EdoneEN{There are exactly two frogs in each row and each column. The frogs decide that 2 of them will jump to a neighbouring empty cell at the same time.
Neighbouring cells have a side in common. After that, there still are exactly 2 frogs in each row and in each column. In how many ways can the frogs do this?}
\def\EdoneLT{Kiekvienoje eilutėje ir kiekviename stulpelyje tupi lygiai dvi varlės (žr. paveikslėlį). Varlės
nusprendė, kad dvi iš jų vienu metu peršoks į gretimus tuščius langelius taip, jog kiekvienoje
eilutėje ir kiekviename stulpelyje tebebus lygiai dvi varlės (gretimi langeliai turi bendrą kraštinę). Keliais būdais
varlės gali tai atlikti?}
\def\EdonePL{W każdym rzędzie i w każdej kolumnie znajdują się dokładnie dwie żaby. Żaby ustaliły, że dwie z nich jednocześnie przeskoczą na
sąsiednie puste pola, ale tak, by nadal w każdym rzędzie i w każdej kolumnie były dokładnie dwie żaby. (Sąsiednie pola to takie, które mają wspólny bok.)
Na ile sposobów żaby mogą to zrobić?}
\def\EdoneRU{На рисунке в каждой строке и каждом столбце сидят ровно две лягушки.
Две из них должны одновременно прыгнуть в соседние пустые клетки так, чтобы в каждой строке и каждом столбце осталось ровно две лягушки
    (соседние клетки имеют общую сторону). Сколькими способами лягушки могут это сделать?}

\def\rEdone{C}
\def\aEdone{\answer{2}{3}{4}{5}{6}}

%E22===============================[======================........]===============================%E22
\def\EdtwoEN{The figure below shows a honeycomb with nine cells. There is honey in some cells.
The number in each cell shows how many neighbouring cells contain honey (neighbouring cells have a side in common).
How many cells contain honey?}
\def\EdtwoLT{Kai kuriose iš devynių korio ląstelių yra medaus (žr. paveikslėlį). Skaičius kiekvienoje ląstelėje rodo,
keliose iš jai gretimų ląstelių yra medaus (dvi gretimos ląstelės turi bendrą kraštinę).
Keliose korio ląstelėse yra medaus?}
\def\EdtwoPL{Plaster jest utworzony z dziewięciu sześciokątnych komórek. W~niektórych z nich jest miód.
Liczba w komórce wskazuje, w ilu sąsiadujących z nią (czyli mających z nią wspólny bok) komórkach jest miód.
W~ilu komórkach znajduje się miód?}
\def\EdtwoRU{Некоторые из девяти ячеек в куске сотов содержат мёд (см.~рисунок).
Число в каждой ячейке показывает, сколько соседних ячеек содержат мёд
    (соседниe  ячейки имеют общую сторону). Сколько ячеек в этом куске сотов содержат мёд?}

\def\rEdtwo{C}
\def\aEdtwo{\answer{4}{5}{6}{7}{8}}

%E23===============================[=======================.......]===============================%E23
\def\EdthreeEN{
Тhree girls go to the tray one after the other and take some cookies.
One of the girls takes all the hearts available on the tray.
One girl takes all the white cookies available on the tray.
One girl takes all the large cookies available on the tray.
However, they do not necessarily take the cookies in this order.
One girl takes $3$ cookies, one takes $6$ cookies and one takes $7$ cookies.
Which of the following sets of cookies does one of these girls take?
    }
\def\EdthreeLT{Ant padėklo buvo didesnių ir mažesnių sausainių (žr.~paveikslėlį). Trys
mergaitės viena po kitos ėjo prie padėklo ir rinkosi sausainius. Viena iš jų pasiėmė visas širdeles, gulėjusias
tuo metu ant padėklo; viena pasiėmė visus baltus sausainius, gulėjusius tuo metu ant padėklo;
viena pasiėmė visus didesniuosius sausainius, gulėjusius ant padėklo (mergaitės rinkosi sausainius nebūtinai
ką tik surašyta tvarka). Viena mergaitė pasiėmė 3 sausainius, viena 6 sausainius, viena 7 sausainius (nežinia
    kuria tvarka). Kurį iš išvardytų rinkinių pasiėmė viena iš mergaičių?}
\def\EdthreePL{Troje dzieci jedno po drugim podchodziło do tacy i brało ciastka, niekoniecznie w przedstawionej kolejności.
Jedno z dzieci wzięło z tacy wszystkie ciastka w kształcie serca. Inne dziecko zabrało wszystkie jasne ciastka.
Jeszcze inne dziecko wzięło wszystkie duże ciastka. Liczba ciastek, które wzięły dzieci, to \mat{3}, \mat{6} i \mat{7} (niekoniecznie w tej kolejności).
Który z poniższych zestawów ciastek wzięło jedno z dzieci?}
\def\EdthreeRU{На подносе лежали печенья различной формы, причём некоторые из них были побольше,
другие – поменьше (см. рисунок). Три девочки поочерёдно подходили к подносу и выбирали себе печенья.
Одна из них взяла себе все сердечки, лежавшие на подносе в тот момент; одна взяла все
белые печенья, лежавшие на подносе в тот момент; одна взяла все большие печенья,
лежавшие на подносе в тот момент (не обязательно в указанном
порядке). Одна девочка взяла 3 печенья, одна 6 печений, одна 7 печений (неизвестно в
каком порядке). Который набор из перечисленных ниже взяла одна из
девочек?}

\def\rEdthree{E}
\def\aEdthree{\granswertwentythree[0.9]{E23-2}{E23-3}{E23-4}{E23-5}{E23-6}}
\def\aEdthreeEXP{\granswerexpecolier[1]{E23-2}{E23-3}{E23-4}{E23-5}{E23-6}}

%E24===============================[========================......]===============================%E24
\def\EdfourEN{There are two types of blocks: white \includegraphics[scale=0.15]{E24-1} and grey \includegraphics[scale=0.15]{E24-2}.
A small cube can be made of 4 white blocks or of 1 white and 1 grey block.  The large cube shown in the picture is made of small cubes.
What is the smallest number of white blocks needed to make this large cube?}
\def\EdfourLT{Stasys turi dviejų rūšių detales: baltas \includegraphics[scale=0.15]{E24-1} ir pilkas \includegraphics[scale=0.15]{E24-2}.
Mažą kubelį galima sudėti arba iš keturių baltųjų detalių, arba iš vienos baltosios ir vienos pilkosios.
Iš mažųjų kubelių Stasys sudėjo didelį kubą (žr. paveikslėlį). Kiek mažiausiai jam reikėjo baltųjų detalių?}
\def\EdfourPL{Sławek ma dwa rodzaje klocków: białe \includegraphics[scale=0.15]{E24-1}
i~szare \includegraphics[scale=0.15]{E24-2}. Mały sześcian buduje albo z czterech białych klocków,
    albo z jednego białego i~jednego szarego klocka. Z małych sześcianów Sławek zbudował większy sześcian --- patrz obrazek.
    Jaka jest najmniejsza liczba białych klocków, których mógł użyć Sławek do swej budowli?}
\def\EdfourRU{Есть два типа деталей: белые \includegraphics[scale=0.15]{E24-1} и серые \includegraphics[scale=0.15]{E24-2}.
Маленький кубик можно составить либо из 4 белых деталей, либо из 1 белой и 1 серой детали.
Показанный на рисунке большой куб составлен из маленьких кубиков.
Какое наименьшее количество белых деталей достаточно для сборки такого куба?}

\def\rEdfour{D}
\def\aEdfour{\answer{8}{11}{13}{14}{23}}

